% 08-privacy.tex: граница приватности, изложенная точно и измеренная там, где её измеряли.
\section{Граница приватности: что верификатору дозволено узнать}
\label{sec:privacy}

Две вещи в Polaris используют нулевое разглашение, одна вещь намеренно не существует, и различие
между ними есть самое важное предложение настоящего раздела. Версия 3 добавляет четвёртое:
противника, написанного, чтобы опровергнуть утверждение, и то, что он нашёл.

\subsection{Три уровня раскрытия}

Всякое событие проверки записывается на одном из трёх уровней. \code{ZERO\_KNOWLEDGE} доказывает,
что действительный токен существует, не записывая, какой именно: идентификатор токена в событии
пуст, а ограничение целостности отвергает всякую строку, утверждающую иное (C2). \code{SELECTIVE}
записывает только названные признаки. \code{FULL} записывает личность, заносится в журнал и
ограничивается по частоте. Уровень решается сервером и не может быть повышен клиентом (C6):
доверяющая сторона не способна обратить проверку с нулевым разглашением в полную, спросив иначе, а
Атлас скрывает местоположение события с нулевым разглашением на сервере. Уровень по умолчанию есть
нулевое разглашение. \sImpl

\subsection{Несвязываемые записи проверок на стороне эмитента}

Первое применение нулевого разглашения есть свойство базы данных эмитента. Проверка в режиме по
умолчанию не хранит идентификатора токена, так что граф проверок, кого где и когда проверяли,
невосстановим из базы данных даже тем, кто держит её целиком. Это структурное утверждение, а не
обещание: репозиторий несёт доказательство сокрытия, показывающее, что скрытое поле невыводимо при
заявленной модели противника, тесты свойств прогоняют каждый путь чтения, а утверждение, что две
доверяющие стороны, объединившие всё, чем владеют, всё равно не свяжут держателя, есть формальная
модель, проверяемая при каждой отправке. Тест свойства, измеряющий риск восстановления, натравливает
противника на изолированные события с нулевым разглашением и утверждает, что его доля успехов не
лучше случайной, а начиная с \code{v9.408} утверждение устроено так, что слабый противник не в
силах польстить результату: две слабости противника были для прежнего устройства невидимы.
Утверждение ограничено точно. Оно покрывает события в режиме нулевого разглашения. Избирательные и
полные события идентификатор токена несут, потому что это и означают те уровни. \sImpl

% Рисунок открывает подраздел под его заголовком, и заголовок с рисунком остаются на одной
% странице: нужное им вместе место резервируется до заголовка, так что страница, которая не
% вмещает обоих, разрывается перед заголовком, а не между ними.
\Needspace{44\baselineskip}
\subsection{Доказательство принадлежности}

\figfile{holder-flow}{Держатель, доверяющая сторона и эмитент, с тем, что каждый видит и чего
каждому видеть не дают. Держатель предъявляет файл; доверяющая сторона решает автономно
отделяемым верификатором или SDK и узнаёт, что действительное удостоверение существует, и его
значение; эмитент записывает, в режиме нулевого разглашения, событие без идентификатора токена и не
держит строки <<кто кого проверил>> для сетевой проверки состояния. Красное ребро есть заявленное
ограничение: при обыкновенном предъявлении значение удостоверения есть устойчивый дескриптор,
одинаковый у всех верификаторов, которым его показывают.}{fig:holder-flow}

Второе применение нулевого разглашения есть объект-доказательство. Дерево Меркла ведёт себя как
криптографический отпечаток целого множества: вместо того чтобы вручать верификатору всё
множество, можно вручить ему короткое доказательство того, что определённый предмет принадлежит
множеству с зафиксированным отпечатком. Polaris фиксирует, для каждой эпохи, корень Меркла над
множеством действующих токенов, а держатель способен произвести на собственном устройстве
доказательство Plonky2 против опубликованного множества анонимности о том, что его токен есть лист
под этим корнем, связанный с эпохой, контекстом и одноразовым числом, не открывая, какой именно
лист. Верификатор видит байты доказательства и открытые входы и восстановить лист не способен; это
свойство системы доказательств, а не правило API. Настройка прозрачна, обряда не требует, а
обязательства построены на хеш-функциях. Глубина дерева по умолчанию, равная четырнадцати,
покрывает заданный схемой потолок эпохи в десять тысяч листьев, так что множество анонимности есть
целая эпоха. На эталонной машине доказательство занимает около 35~мс на изготовление и около
13~мс на проверку, а вердикт удостоверяется двумя свидетелями на уровне утверждения.
Доказательство несёт поверификаторный нуллификатор,
\code{Poseidon(секрет\,||\,область\,||\,эпоха)}, так что доверяющая сторона способна соблюсти
правило \emph{один человек однажды} в собственной области, тогда как другая её нуллификаторы
соотнести не сумеет. \sImpl

Чего доказательство не делает, изложено с той же тщательностью. Оно доказывает принадлежность и
привязку, и больше ничего: действующий ли токен, дозволен ли он для контекста и не отозван ли,
решается при составлении эпохи, а не внутри схемы. Одноразовое число мешает подставить
перехваченное доказательство в другой контекст; в сети хранилище одноразовых чисел отвергает
переигранный набор, а автономно верификатор состояния не держит, так что доверяющая сторона,
которой нужно одноразовое использование, выдаёт вызов. Нуллификатор не прячет держателя от
\emph{эмитента}, который выводит каждый лист, чтобы построить дерево, и способен вычислить
нуллификатор любого участника в любой области; свойство действует между доверяющими сторонами. А
схема поверх библиотеки доказательств внешнего криптографического аудита не проходила; книга учёта
обоснованности репозитория говорит, что защищать ею настоящие личности в поставляемом виде не
следует. Сама библиотека доказательств была заново оценена против своей преемницы на
\code{v9.361} и оставлена, причём поводы, которые вновь откроют этот вопрос, записаны. \sExt{} для
аудита.

\subsection{Ограниченность, а не постоянство, и что ограниченность оказалась означать}

Polaris есть система предъявления полного удостоверения, и до \code{v9.353} следствие было
плоским: субъектом входного токена был хеш удостоверения, одно и то же значение у всякой
доверяющей стороны, так что две из них могли соединить свои таблицы пользователей в точности и
навсегда. Субъект и дескриптор предъявления теперь выводятся в собственной области доверяющей
стороны, устойчивы там, где учётной записи это нужно, и неузнаваемы у следующего верификатора;
поскольку субъект есть то значение, которое доверяющая сторона \emph{записывает}, а записанные
значения суть те, что объединяют, продают, истребуют по повестке и теряют при утечках, это есть
половина, которая важнее всего на деле. Другая половина не изменилась: обыкновенное предъявление
по-прежнему показывает верификатору устойчивое \code{token\_value}, подпись эмитента и открытый
ключ держателя, так что гарантия есть гарантия о том, что верификатору следует \emph{хранить}, а
не о том, что ему \emph{показывают}, и верификатор сообщает, что именно он выдаёт:
\code{exposed} для обыкновенного предъявления, \code{bounded} для пути с нулевым разглашением.
\sImpl

Версия 3 способна сказать, что \code{bounded} оказалось означать по измерению, потому что по
операционному договору лаборатория репозитория задала вопрос прямо: имея полные протоколы двух
предъявлений, какое преимущество есть у сговорившейся пары верификаторов при решении вопроса
\emph{тот же ли это держатель}? Три находки, в порядке, в котором они были сделаны.

\begin{description}
  \item[Вердикт означал не то, что говорил.] \code{bounded} сообщалось всякий раз, когда
    присутствовал нуллификатор с областью. Он ни разу не проверял посылки, изложенной в его
    собственной документации. Предъявление, несущее нуллификатор \emph{и} устойчивое значение
    токена, ключ эмитента и подпись, объявлялось ограниченным, тогда как два верификатора могли
    связать его одним сравнением строк. Исправлено в тот же день: \code{bounded} теперь требует
    нуллификатора и отсутствия всякого поля, одинакового у разных верификаторов, включая открытый
    ключ держателя, который устойчив по построению и ровно поэтому попарный дескриптор из него
    выводится, а не показывается. Исправление поправило вердикт; ни одно предъявление приватнее,
    чем было, оно не сделало.
  \item[Никакого преимущества сверх множества анонимности, и ничто не задаёт ему нижнего предела.]
    Противник, который уплощает всякую величину в протоколе, отбрасывает постоянное по всему
    населению и сопоставляет по остальному, имея положительный контроль, связывающий открытое
    население с точностью единица, достигает на ограниченном населении точности, следующей за
    $1/N$ при всяком размере от двух держателей до двухсот. Внутри эпохи привязка к области делает
    то, что обещает слово, против этого противника. А $N$ есть число участников эпохи. Схема
    стесняет зафиксированное число эпохи промежутком от одного до десяти тысяч, и ничто не задаёт
    ему нижнего предела выше единицы, так что эпоха, закрывшаяся тремя участниками, даёт
    сговорившейся паре догадку один к трём, а закрывшаяся единственным участником опознаёт
    держателя прямо, тогда как верификатор по-прежнему, и правильно, сообщает протокол как
    ограниченный. Число читается с подписанной контрольной точки; ничто не обязывает доверяющую
    сторону смотреть, и ни один вердикт о нём не упоминает. Сообщать его было бы новой продуктовой
    гарантией, создавать которые лабораторной работе не дано, поэтому это записано как ограничение
    для развёртывающей организации.
  \item[Чего противник не видит.] Он сопоставляет по точному равенству и ни по чему иному.
    Покажи ему население, идеально связываемое, но без единого равного поля, и он останется на
    уровне случайности. Значит, нулевой результат означает, что ни одно поле не \emph{одинаково} у
    разных верификаторов, а не что ни одно не \emph{коррелировано}. Единственный канал, который он
    способен прочесть, есть длина протокола, и она несла весь остаток выше случайности в том
    контроле; ничто в ограниченном предъявлении Polaris сегодня не меняет размера от держателя к
    держателю, а развёртывание, в котором наборы раскрываемых признаков или необязательные
    элементы заставили бы его меняться, этот канал откроет. Размер предъявления, временные
    характеристики, метаданные эмитента, предметы о состоянии и структура протокола не измерены, и
    лаборатория говорит об этом, а не экстраполирует.
\end{description}

Путь, которым в действительности воспользовался кошелёк постороннего, режет в другую сторону.
Предъявление SD-JWT VC поверх OpenID4VP, измеренное тем же способом, оставляет девять значений
одинаковыми у двух верификаторов: ключ держателя, подпись эмитента, дайджесты раскрытий, хеш
предъявления и его длину, и все они подписаны эмитентом, так что держатель изменить их не может.
Это известное свойство многократного предъявления одного такого удостоверения; смягчение, принятое
в этой экосистеме, есть пакетная выдача, много удостоверений по одному ключу держателя на каждое,
чего репозиторий не реализует. Принятие формата экосистемы внесло связываемость этой экосистемы;
до измерения это не было видно, и это не повод прекращать взаимодействие. Итак, межверификаторная
соотносимость ограничена на родном пути с нулевым разглашением, в пределах измеренного выше, и
открыта на пути совместимости. \emph{Несвязываемость по умолчанию} по-прежнему означает записи
эмитента в режиме нулевого разглашения. \sImpl{} для выводов и измерения, с остатком, изложенным,
а не закрытым.

\subsection{Ограниченные поверхности}

Оставшиеся механизмы приватности ограничивают то, что поверхность способна вернуть, а не то, что
хранится. Всякая карта и всякая сводка API имеют жёсткий потолок (C8), так что всё население нельзя
извлечь через операторский атлас, а запрос на шесть с половиной миллионов ячеек кластера
возвращает самое большее пять тысяч. Страницы расследования по одной сущности построены так, чтобы
отказывать в анализе связей между людьми, и они подлежат аудиту: таблица мета-аудита записывает,
кто читал какую поверхность аудита, с каким фильтром и сколько строк вернулось, и она работает
только на добавление. Самое навязчивое полномочие органа, вытянуть всю историю проверок одного
названного человека по ордеру, не оставляло записи о собственном применении всё время своего
существования, потому что чтение сидело за хранимой процедурой и ничто, искавшее \code{SELECT},
его не находило; начиная с \code{v9.382} оно пишет ту же строку, а структурное правило
распространяет это на всякое чтение, спрятанное за процедурой или представлением. Слою структуры
органов, Афине, пять проверок запрещают ссылаться на любую таблицу или поле физического лица
(раздел~\ref{sec:cognition}). А отметки времени, служба отметок времени и шлюз обмена построены
так, чтобы ничего не узнавать о том, чего касаются, что описывает раздел~\ref{sec:institutions}.
\sImpl

\begin{layercard}{Карточка слоя: граница приватности}
\qa{Что это}{Три уровня раскрытия с ограничением CHECK за уровнем по умолчанию; закрепление на сервере; доказательство сокрытия и его противник; доказательство принадлежности эпохе, его нуллификатор и его второй свидетель; поверификаторный дескриптор; ограниченные сводки; поднадзорная поверхность расследования по одной сущности.}
\qa{Зачем оно существует}{Журнал проверок есть хребет слежки, если он структурно не лишён способности им быть.}
\qa{Чему доверяет}{Ограничению C2; обоснованности библиотеки доказательств; модели противника в доказательстве сокрытия и заявленным слепым пятнам противника по связываемости.}
\qa{Чему отказывает}{Утверждению клиента о собственном уровне раскрытия; идентификатору токена в строке с нулевым разглашением; неограниченному чтению; узлу-человеку в слое органов; вердикту \code{bounded} над протоколом, в котором хоть одно поле одинаково у разных верификаторов.}
\qa{Что видит}{При нулевом разглашении: что действительный токен был проверен в такое-то время в таком-то контексте. При полном: личность, с занесением в журнал.}
\qa{Чего не видит}{Какой именно токен, в режиме нулевого разглашения, даже из всей базы данных; и размер толпы, в которой прячется данное доказательство, о чём не сообщает ни один вердикт.}
\qa{Если откажет}{Идентификатор в строке с нулевым разглашением нехраним; убранный потолок роняет сборку; ссылка на человека в Афине роняет пять проверок; противник, ослабленный ради лестного результата, роняет собственный контроль набора тестов свойств. Межверификаторная соотносимость при обыкновенном предъявлении есть не отказ, а описанная граница.}
\qa{Кем проверяется}{CHECK C2 и модель C2, \code{check\_c6\_atlas\_redacts\_zk\_location}, набор тестов свойств сокрытия, \code{check\_redaction\_adversary\_is\_not\_flattered}, \code{check\_pairwise\_presentation}, \code{check\_scoped\_nullifier}, \code{check\_audited\_reads\_are\_logged}, дифференциал второго свидетеля ZK, \code{check\_athena\_no\_person}.}
\qa{Сейчас или потом}{Сейчас. Предъявление, устойчивое к связыванию, на пути совместимости и нижний предел размера толпы суть другая система и потом.}
\qa{Внутри и снаружи}{Внутри: состояние. Снаружи: держатели и доверяющие стороны, пользующиеся удостоверением.}
\end{layercard}

\nextlayer{Удостоверение, ограниченное в том, что открывает, всё равно нужно использовать: держать,
предъявлять, принимать и входить по нему. Следующий слой есть место, где удостоверение встречает
программы, которые не принадлежат эмитенту, и где, впервые, эти программы не принадлежат и автору.}
